\section{Background}\label{sec:background}

\subsection{Key Concepts}

\subsubsection{Formula 1 Tire Degradation and Pit Stop Strategy}

Formula 1 regulations require each driver to pit atleast once, and use at least two different tire compounds during a dry race. The three primary 
compounds (SOFT, MEDIUM, and HARD) differ in grip level and durability. SOFT tires offer peak grip and are considered the fastest, however, only
for $\sim$15--25 laps before 
significant degradation; HARD tires are considered the slowest but last $\sim$35--50 laps and generate less grip throughout. Furthermore, tire degradation is influenced by numerous factors such as track temperature, compound type, driving style, etc.
As a stint progresses, lap times increase monotonically due to tire wear, providing a time-series signal 
that, in theory, should predict when pitting becomes beneficial.

However, the pit stop decision is not just about when tires degrade. It involves strategic considerations such as the undercut 
(pitting early to gain position through fresher tires), the overcut (staying out while a competitor pits, maintaining track 
position), and the safety car window (pitting under a safety car to take a ``free'' stop). This project addresses the tire 
degradation component only; competitor position and safety car timing are treated as out-of-scope.

\subsubsection{Recurrent Neural Networks for Time Series}

Long Short-Term Memory (LSTM) networks \cite{hochreiter1997} and Gated Recurrent Units (GRUs) \cite{cho2014} are the dominant 
architectures for sequential prediction tasks. Both maintain a hidden state that summarizes prior context, updated at each
 timestep through gating mechanisms.

LSTM cells use three gates (input, forget, output) and a separate cell state, giving them the concept of `memory', and strong capacity for long-range
 dependencies at the cost of higher parameter counts. GRU cells use two gates (reset, update) and no separate cell state,
  making them faster to train and less prone to overfitting on moderate-sized datasets. For short sequences (10--20 laps),
   the memory advantage of LSTM is less relevant, and GRU has been shown to match or outperform LSTM in similar settings\cite{chung2014}.

Attention mechanisms\cite{bahdanau2015} extend recurrent models by computing a weighted sum over 
all hidden states rather than using only the final state, allowing the model to selectively focus on the most 
informative timesteps. Their benefit is most pronounced on sequences of 20 or more steps where individual timesteps 
carry distinct signals.

\subsubsection{Delta Prediction for Degradation Modeling}

Standard sequence models trained to predict absolute lap time face a systematic bias: the model must first ``explain'' 
the circuit's baseline lap time before it can model the degradation delta on top of it. This makes it difficult to learn 
compound-specific degradation slopes.

Reformulating the target as $\Delta t = t_{n+1} - t_n$ (lap-to-lap change) decouples the degradation signal from the baseline, 
allowing the model to focus exclusively on the rate of change. This formulation is directly motivated by degradation physics, which suggests that
tire wear increases lap time at a rate that depends on compound hardness and track temperature, not on the absolute lap time.

\subsection{Related Work}

\subsubsection{Race Strategy Simulation}

Commercial F1 strategy tools (e.g., McLaren Electronic Systems, Williams Racing Analytics) are widely reported to use physics-based tire models combined with Monte Carlo simulations over opponent strategies. These systems require real-time telemetry and proprietary tire compound data, neither of which are available for public use.

Academic work on F1 strategy is severely limited due to the lack of publicly available data. Heilmeier et al.~\cite{heilmeier2020} develop a race simulation framework for circuit motorsports that incorporates tire degradation, fuel load, and safety car probability, using a parametric tire model fitted on historical timing data. Their approach requires compound-specific parameters typically sourced from team data and focuses on pre-race strategy planning rather than real-time window detection.

Tulabandhula and Rudin~\cite{tulabandhula2014} apply machine learning to sports decision problems more broadly, demonstrating that data-driven approaches can approximate expert judgment in complex sequential settings. Reinforcement learning has also been informally explored for racing strategy, but these approaches require simulator environments that are difficult to calibrate without access to official team data.

\subsubsection{Tire Degradation Modeling}

Tire degradation has been modeled as a parametric function of stint length in physics-based simulators~\cite{heilmeier2020}. Data-driven approaches in the broader motorsport literature include regression over lap time residuals and recurrent models trained on historical sector times. To my knowledge, no published work applies compound-specific GRU models to pit stop window detection using publicly available lap data.

\subsubsection{Publicly Available F1 Data}

The FastF1 Python library \cite{fastf1} provides a clean interface to the official Formula 1 Livetiming API, including lap times, tire compound, weather, and driver information for all races since 2018. It is the only fully open data source for F1 lap timing and is used in this project as the sole data source.
